\documentclass[10pt]{article}
\usepackage[accepted]{icml2026}

\usepackage[T1]{fontenc}
\usepackage[utf8]{inputenc}
\usepackage{lmodern}
\usepackage{microtype}
\usepackage{amsmath, amssymb, amsthm, mathtools}
\usepackage{booktabs}
\usepackage{array}
\usepackage{enumitem}
\usepackage[numbers,sort&compress]{natbib}
\usepackage{xcolor}
\usepackage{hyperref}
\usepackage{float}

\hypersetup{
  colorlinks=true,
  linkcolor=blue!50!black,
  citecolor=blue!50!black,
  urlcolor=blue!50!black
}

\newtheorem{theorem}{Theorem}
\newtheorem{corollary}{Corollary}

\newcommand{\E}{\mathbb{E}}
\newcommand{\KL}{D_{\mathrm{KL}}}
\newcommand{\Var}{\mathrm{Var}}
\newcommand{\diag}{\mathrm{diag}}
\newcommand{\HM}{\mathrm{HM}}
\newcommand{\AM}{\mathrm{AM}}
\newcommand{\meta}{\mathrm{meta}}
\newcommand{\LoRA}{\mathrm{LoRA}}
\newcommand{\RESULTPENDING}{\textcolor{red!70!black}{\textbf{RESULT PENDING}}}

\icmltitlerunning{Meta-SWAG}

\title{\textbf{Meta-SWAG: Bayesian Posterior Approximation over Markovian Learning Trajectories for Multi-Agent Policy Gradient and LLM Alignment}}
\author{}
\date{}

\begin{document}

\twocolumn[
\maketitle
\begin{center}
\vspace{-1.75em}
\begin{minipage}{0.92\textwidth}
\begin{abstract}
Posterior approximation is now a standard tool for uncertainty in supervised deep learning, but comparable machinery is still missing in settings where the training trajectory is generated by an evolving learning process rather than a fixed objective. We introduce \textbf{Meta-SWAG}, a SWAG-style Gaussian posterior approximation over checkpoints drawn from a \textbf{Markovian learning trajectory}. Our starting point is that multi-agent policy gradient and LLM alignment both produce adaptive checkpoint sequences whose effective objective shifts over time, even though the source of that non-stationarity differs across the two domains. Markovianity alone does not justify a Gaussian approximation; the approximation becomes meaningful only after retaining checkpoints inside a stable, mixing region of the trajectory. Under that stronger regime, we define a posterior $q_T(\theta)=\mathcal{N}(\mu_T,\Sigma_T)$ over trainable parameters, restricted to LoRA adapters in the LLM setting. We present two formal results and one conceptual limitation result: under explicit evidence-variance alignment assumptions, Bayesian model averaging improves over point estimates through a variance reduction term; the posterior covariance tracks local trajectory geometry near stable regions; and posterior quality is bottlenecked by a self-knowledge gap between an updater and the policy model it relies on. We also introduce Goodhart-resilient weighting rules that prevent posterior collapse onto proxy-maximizing checkpoints. Empirically, the current draft pairs completed matrix and iterated-game scaffolds with a LoRA-based DPO/PPO/GRPO protocol that remains partially pending.
\end{abstract}
\end{minipage}
\end{center}
\vspace{0.8em}
]

\section{Introduction}

Uncertainty estimates for learned policies remain underdeveloped in exactly the domains where they matter most. In multi-agent reinforcement learning, learning dynamics are shaped by the simultaneous adaptation of other agents \citep{kim2021meta,giannou2022convergence,foerster2018lola}; in LLM alignment, the object being optimized is a changing policy produced by DPO, PPO, or GRPO updates rather than a fixed supervised predictor \citep{rafailov2023dpo,schulman2017ppo,shao2024deepseekmath}. Both settings offer local convergence or optimization stories, but neither comes with a standard posterior approximation that can be used for model averaging, geometry diagnostics, or uncertainty-aware deployment.

SWAG and related subspace methods are attractive candidates because they give cheap Gaussian posteriors from training trajectories \citep{maddox2019swag}. The obstacle is that vanilla SWAG presumes a stationary loss landscape and approximately exchangeable late iterates. Those assumptions fail in raw MARL and raw fine-tuning trajectories because the effective objective changes as the environment, the peer policies, or the aligned policy itself changes. Our central claim is therefore narrow: a SWAG-style posterior can still be used on \textbf{Markovian learning trajectories} once checkpoint retention is restricted to a stable, mixing region in which local moment estimation is meaningful. The LLM case is therefore not a claim that the optimizer is an agent; it is the milder claim that fine-tuning induces a checkpoint process whose states evolve Markovianly under repeated alignment updates.

This paper introduces Meta-SWAG, a Gaussian posterior approximation over adaptive checkpoint sequences. The method is built around one clean object, $q_T(\theta)=\mathcal{N}(\mu_T,\Sigma_T)$, with the posterior restricted to LoRA adapter parameters in the LLM setting. We make three claims. First, under an explicit evidence-variance alignment assumption, Bayesian model averaging over the posterior reduces predictive variance relative to nearby point estimates and uniform ensembles. Second, the posterior covariance captures local trajectory geometry, so uncertainty structure becomes a diagnostic rather than just an interval estimate. Third, posterior quality is constrained by the quality of the updater's internal policy model; in the alignment setting this is presented as a conceptual limitation principle rather than a completed empirical result. The paper does not claim that MARL and LLM alignment share identical non-stationarity; it claims that the same retained-checkpoint posterior construction applies in both domains once their respective stable-region assumptions are made explicit.

Empirically, we design the paper around two validations. The first uses matrix games and Kim-style iterated games to test the variance and geometry claims in a controlled setting. The second targets LoRA-based alignment of a Llama-family model under DPO, PPO, and GRPO, with reward-overoptimization measured by a gold reward model and best-of-$n$ sampling \citep{gao2023overoptimization,rafailov2024daa}. The current draft keeps the full LLM protocol and figure structure even where final runs are still pending, so the contribution should be read as a sharpened method paper with completed small-scale evidence rather than as a fully finished large-scale empirical study.

\section{Formal Setup}

\subsection{Adaptive Checkpoint Process}

Let $\Theta \subseteq \mathbb{R}^d$ denote the parameter space of the trainable variables. In MARL, $\theta_t \in \Theta$ denotes the joint policy parameters of all agents; in LLM alignment, $\theta_t$ denotes the trainable LoRA adapter parameters with the base model frozen. All formal claims are stated for the chosen trainable coordinate space $\Theta$; in the LLM experiments this means the claims are about posteriorization in LoRA space, not about full-model uncertainty. We study a retained checkpoint process
\begin{equation}
\theta_{t+1} = \mathcal{T}(\theta_t,\xi_{t+1}),
\end{equation}
where $\mathcal{T}$ is a measurable update operator and $\xi_{t+1}$ collects the randomness induced by minibatch sampling, rollout noise, and environment stochasticity. For the paper's theory, the object of interest is the post-burn-in process $\{\theta_t\}_{t \geq t_0}$ rather than a single terminal iterate.

This formulation is intentionally mild. It does not require the optimizer itself to be an agent. It only requires that repeated updates induce a Markovian process on checkpoints. In multi-agent policy gradient, the effective objective shifts because peer adaptation changes the learning environment \citep{kim2021meta,giannou2022convergence}. In LLM alignment, the effective objective shifts because repeated policy updates change the policy whose outputs define the next gradient step \citep{rafailov2023dpo,schulman2017ppo,shao2024deepseekmath}.

\subsection{Retention Rule and Evidence Scores}

Let $\mathcal{K}_T = \{t_1,\dots,t_K\} \subseteq \{t_0,\dots,T\}$ denote the retained checkpoint indices after burn-in, ordered increasingly, and write $\vartheta_k := \theta_{t_k}$ for the retained checkpoints. Each retained checkpoint is assigned an evidence score $m_k \in \mathbb{R}$ computed on held-out data:
\begin{itemize}[leftmargin=1.5em]
  \item in MARL, $m_k$ can be a held-out meta-value or adaptation score;
  \item in LLM alignment, $m_k$ can be a held-out preference accuracy, reward-model score, or other validation statistic.
\end{itemize}
The weighting map is a measurable function $\psi : \mathbb{R} \to \mathbb{R}_{>0}$, producing normalized weights
\begin{equation}
w_k = \frac{\psi(m_k)}{\sum_{j=1}^K \psi(m_j)}.
\end{equation}
The naive choice is $\psi(m)=e^{\beta m}$; the Goodhart-resilient variants introduced later modify $\psi$ or tune $\beta$ to avoid posterior collapse.

\subsection{Standing Assumptions}

We will invoke the following standing assumptions, with theorem-specific strengthening stated where needed.
\begin{enumerate}[leftmargin=1.5em, itemsep=0.25em]
  \item \textbf{(A1) Markov checkpoint process.} The post-burn-in checkpoint process is time-homogeneous and Markov on $\Theta$.
  \item \textbf{(A2) Stable-region retention.} The retained checkpoints lie in a compact region $\mathcal{R} \subseteq \Theta$ with finite second and fourth moments.
  \item \textbf{(A3) Mixing.} The retained process is geometrically mixing on $\mathcal{R}$, so empirical first and second moments obey the usual concentration and law-of-large-numbers behavior.
  \item \textbf{(A4) Local linearization.} On $\mathcal{R}$, the update process admits a first-order linearization around a stable reference point or recurrent orbit, with square-integrable residual noise.
  \item \textbf{(A5) Evidence regularity.} The score map and weighting rule produce strictly positive normalized weights; additional alignment between evidence and checkpoint quality is imposed only in Theorem~1.
\end{enumerate}

\subsection{Why SWAG Needs a New Object}

Standard SWAG approximates the posterior over parameters by fitting a Gaussian to late training iterates \citep{maddox2019swag}. The approximation is compelling when the iterates come from a relatively stationary neighborhood of a single basin, so that first and second moments summarize the local geometry well. In raw RL or raw fine-tuning, however, these assumptions fail for two reasons. First, the effective objective is non-stationary because the data distribution changes with the policy. Second, late iterates need not be exchangeable when the update rule is itself part of the object being studied.

The Markovian checkpoint formulation identifies the right stochastic object, but it does not by itself restore the Gaussian approximation. The additional work is done by retention inside a stable region together with mixing and local linearization assumptions, which let late-trajectory moments summarize one coherent part of the adaptive process. Meta-SWAG therefore approximates uncertainty over adaptive policies only after this restricted late-trajectory regime is isolated; it is not a claim that arbitrary non-stationary training traces are Gaussian.

\section{Meta-SWAG}

\subsection{Posterior Definition}

Given retained checkpoints $\{\vartheta_k\}_{k=1}^K$ and normalized weights $\{w_k\}_{k=1}^K$, Meta-SWAG defines a Gaussian posterior
\begin{equation}
q_T(\theta)=\mathcal{N}(\mu_T,\Sigma_T),
\end{equation}
with weighted mean
\begin{equation}
\mu_T = \sum_{k=1}^K w_k \vartheta_k,
\end{equation}
diagonal second-moment accumulator
\begin{equation}
v_T = \sum_{k=1}^K w_k (\vartheta_k-\mu_T)\odot(\vartheta_k-\mu_T),
\end{equation}
and low-rank deviation matrix
\begin{equation}
D_T = \left[\sqrt{w_k}(\vartheta_k-\mu_T)\right]_{k \in \mathcal{J}_T},
\end{equation}
where $\mathcal{J}_T \subseteq \{1,\dots,K\}$ indexes the retained low-rank buffer. The covariance estimate is
\begin{equation}
\Sigma_T = \frac{1}{2}\diag(v_T) + \frac{1}{2(|\mathcal{J}_T|-1)}D_TD_T^\top.
\end{equation}

For LLM alignment, all quantities are computed in LoRA space only. This restriction is both computationally necessary and conceptually appropriate: the posterior lives on the subspace actually updated by alignment. The scope of theorems in the LLM setting should therefore be read accordingly. They characterize uncertainty and geometry in the trainable adapter space, not in the frozen base-model directions.

\subsection{Goodhart-Resilient Weighting}

The baseline softmax rule
\begin{equation}
w_k \propto e^{\beta m_k}
\end{equation}
can collapse onto a single proxy-maximizing checkpoint. To prevent this, we use Goodhart-resilient weighting rules.
\begin{enumerate}[leftmargin=1.5em, itemsep=0.25em]
  \item \textbf{ESS-constrained softmax.} Choose $\beta$ so that
  \begin{equation}
  ESS(w) := \frac{(\sum_k w_k)^2}{\sum_k w_k^2}
  \end{equation}
  stays above a prescribed floor.
  \item \textbf{Thresholded satisficing.} Assign uniform weight to checkpoints whose evidence exceeds a validation threshold and zero weight otherwise.
\end{enumerate}
These rules keep the method inside the same Gaussian family while making the anti-collapse constraint explicit.

\subsection{Bayesian Alignment Landscape}

Recent work has brought Bayesian ideas into LLMs through local curvature approximations, adapter-space posteriors, and Bayesian reward models \citep{yang2024laplacelora,yang2024brm}. Meta-SWAG complements rather than replaces these approaches. Laplace-LoRA approximates posterior uncertainty around one adapted solution using local curvature in LoRA space; Bayesian reward-model approaches place uncertainty on the learned reward function; Meta-SWAG instead places uncertainty on the aligned policy trajectory itself.

\begin{table}[H]
\centering
\small
\begin{tabular}{>{\raggedright\arraybackslash}p{0.22\linewidth} >{\raggedright\arraybackslash}p{0.24\linewidth} >{\centering\arraybackslash}p{0.15\linewidth} >{\centering\arraybackslash}p{0.15\linewidth} >{\centering\arraybackslash}p{0.17\linewidth}}
\toprule
Method & Posterior object & Extra models & Trajectory-aware & DPO/PPO/GRPO \\
\midrule
Laplace-LoRA & Local adapter posterior around one solution & No & No & Usually yes \\
Bayesian reward modeling & Reward model posterior & Yes & Indirectly & Yes \\
\textbf{Meta-SWAG} & \textbf{Posterior over aligned policy checkpoints} & \textbf{No} & \textbf{Yes} & \textbf{Yes} \\
\bottomrule
\end{tabular}
\caption{Positioning Meta-SWAG against nearby Bayesian alignment methods.}
\end{table}

The distinction matters because reward uncertainty and policy uncertainty answer different questions. Reward-model posteriors quantify uncertainty about the evaluator. Meta-SWAG quantifies uncertainty about the policy that will actually be sampled at deployment time.

\begin{figure}[H]
\centering
\fbox{
\begin{minipage}{0.94\linewidth}
\small
\textbf{Figure 1 (placeholder): Domain Bridge for Meta-SWAG.}

\vspace{0.5em}
\textbf{Left:} multi-agent policy gradient as a Markovian learning trajectory over joint policy checkpoints.

\textbf{Center:} both domains feed the same Goodhart-aware Gaussian posterior $q_T(\theta)$.

\textbf{Right:} LLM alignment as a Markovian learning trajectory over LoRA checkpoints, where updates are DPO/PPO/GRPO steps.
\end{minipage}
}
\caption{Shared retained-checkpoint abstraction: the same posterior construction applies to two adaptive Markovian checkpoint processes once domain-specific stable-region assumptions are imposed.}
\end{figure}

\subsection{Algorithm}

\begin{figure}[H]
\centering
\fbox{
\begin{minipage}{0.94\linewidth}
\small
\textbf{Algorithm 1: Meta-SWAG}
\begin{enumerate}[leftmargin=1.5em, itemsep=0.2em, topsep=0.4em]
  \item Inputs: retained checkpoints $\{\vartheta_k\}_{k=1}^K$, evidence scores $\{m_k\}_{k=1}^K$, weighting rule $\psi$, low-rank index set $\mathcal{J}_T$, sample count $S$.
  \item Compute normalized weights $\{w_k\}_{k=1}^K$ using either ESS-constrained softmax weighting or thresholded satisficing weights.
  \item Accumulate weighted mean $\mu_T = \sum_k w_k \vartheta_k$.
  \item Accumulate diagonal moment $v_T = \sum_k w_k (\vartheta_k-\mu_T)\odot(\vartheta_k-\mu_T)$.
  \item Form low-rank deviation matrix $D_T$ from the retained weighted deviations indexed by $\mathcal{J}_T$.
  \item Set $\Sigma_T = \frac{1}{2}\diag(v_T) + \frac{1}{2(|\mathcal{J}_T|-1)}D_TD_T^\top$.
  \item Draw posterior samples $\theta^{(s)} \sim \mathcal{N}(\mu_T,\Sigma_T)$ for $s=1,\dots,S$.
  \item Average policy outputs or downstream scores across posterior samples.
\end{enumerate}
\vspace{0.3em}
\textbf{Outputs:} posterior $q_T(\theta)=\mathcal{N}(\mu_T,\Sigma_T)$, posterior samples, Bayesian model average.
\end{minipage}
}
\end{figure}

The training-time cost is the same order as SWAG: one running mean, one diagonal accumulator, and a bounded low-rank buffer. The inference-time cost is $S$ forward passes or policy evaluations. In the LLM setting, the dominant practical choice is to sample LoRA adapters and merge them into the frozen base model at evaluation time.

\subsection{Design Choices}

\paragraph{Why evidence-weight the mean?}
A plain average treats all checkpoints as equally informative, which is undesirable when the meta-objective clearly prefers some stages of training. Evidence weighting makes the posterior predictive closer to the checkpoints that score well on held-out meta-value or alignment criteria without collapsing fully to a point estimate.

\paragraph{Why guard against Goodhart collapse?}
If $m_t$ is itself a proxy objective, naive weighting can over-concentrate posterior mass on the same over-optimized checkpoint that a point estimate would select. ESS-constrained weighting makes this failure mode explicit and prevents collapse by guaranteeing posterior support over multiple retained checkpoints. Thresholded satisficing provides a simpler alternative that preserves variation across all ``good enough'' models.

\paragraph{Why keep a low-rank covariance term?}
The diagonal term captures marginal uncertainty, but the low-rank factor is what carries basin geometry. The second theorem will use this structure directly: if the trajectory stays in a stable basin, the dominant eigenvectors of $\Sigma_T$ track the locally soft and stiff directions of the meta-objective.

\paragraph{Why a joint posterior rather than per-agent posteriors?}
In MARL the relevant object is the joint policy, not a bag of isolated marginals. In alignment the same logic reappears as base-policy and aligner coupling: what matters is uncertainty over the adapted policy that results from the update process, not uncertainty over individual training steps in isolation.

\section{Theory}

\subsection{Variance Reduction Through Bayesian Model Averaging}

\paragraph{Assumption: Evidence-variance alignment.}
For the retained checkpoints, suppose the weighting rule is monotone in an evidence score $m_t$ and that higher evidence corresponds to lower checkpoint-level predictive variance, i.e., $m_t \ge m_s \Rightarrow V_t \le V_s$. This is an idealized assumption, not a generic property of arbitrary validation metrics.

\begin{theorem}[Idealized variance reduction under evidence-variance alignment]
Let $q_T$ be the Meta-SWAG posterior constructed from retained checkpoints satisfying Assumptions (A1)--(A5), and let $\bar{\pi}_{q_T}$ denote the posterior predictive policy. Under evidence-variance alignment, the within-checkpoint component of the predictive variance of $\bar{\pi}_{q_T}$ is bounded by the corresponding uniformly weighted checkpoint average up to the harmonic-to-arithmetic mean ratio of checkpoint variances:
\begin{equation}
\Var_{\mathrm{within}}(\bar{\pi}_{q_T})
\leq
\frac{\HM(V_1,\ldots,V_K)}{\AM(V_1,\ldots,V_K)}
\Var_{\mathrm{within}}(\pi_{\mathrm{unif}})
\leq
\Var_{\mathrm{within}}(\pi_{\mathrm{unif}}),
\end{equation}
where $V_k$ denotes checkpoint-level predictive variance and $\pi_{\mathrm{unif}}$ denotes the uniformly weighted checkpoint ensemble. If the between-checkpoint term remains controlled inside one stable region, this also yields a reduction relative to nearby point estimates.
\end{theorem}

\paragraph{Proof sketch.}
The predictive variance of the Bayesian model average decomposes into a within-checkpoint term and a between-checkpoint term. Under evidence-variance alignment, monotone evidence weighting assigns larger mass to lower-variance checkpoints, so the within-checkpoint term is controlled by the harmonic mean rather than the arithmetic mean. The argument mirrors the AM-HM improvement from evidence-weighted policy gradients, but the random object is now the posterior predictive policy rather than the gradient estimator. Extending the comparison from a uniformly weighted ensemble to a single point estimate requires the additional single-region control stated above.

\begin{corollary}[Alignment]
Assume in addition that, on the retained LoRA posterior support, the best-of-$n$ proxy-gold reward gap is locally upper-bounded by a monotone function of a predictive-dispersion functional that is itself controlled by the within-checkpoint predictive variance. Then any Meta-SWAG weighting rule that strictly reduces the within-checkpoint predictive variance also weakly reduces that local upper bound on reward overoptimization relative to the corresponding point-estimate or uniform-ensemble baseline.
\end{corollary}

\paragraph{Significance.}
The theorem itself is only a variance statement. The alignment corollary is useful only after adding an explicit bridge from predictive dispersion to proxy-gold reward gap, which is exactly the step the empirical protocol must test rather than assume.

\subsection{Posterior Covariance as Basin Geometry}

\begin{theorem}[Posterior covariance identifies local trajectory geometry]
Suppose the retained trajectory remains in a stable local region of the adaptive optimization process, the dynamics admit a local linearization in that region, and the retained checkpoint process is geometrically mixing. Then, as $T \to \infty$, the Meta-SWAG covariance $\Sigma_T$ converges to a matrix whose leading eigenspaces agree with the dominant weakly-curved directions of the local trajectory geometry, while concentration in sharper directions scales with the local second-order structure of the effective objective.
\end{theorem}

\paragraph{Proof sketch.}
Near a stable local region, the adaptive dynamics admit a local linearization whose covariance is governed by the Jacobian and Hessian of the effective objective along the trajectory. The diagonal-plus-low-rank estimator is consistent for the dominant second-moment structure of this local process. Geometric mixing ensures that the checkpoint buffer estimates local covariance rather than transient drift between unrelated regions. The leading eigenvectors therefore align with soft directions of the local orbit or basin, while narrow directions shrink under repeated averaging.

\begin{corollary}[Conditional diagnostic interpretation under subspace-sensitive shift]
Assume additionally that deployment shift perturbs model behavior primarily through the retained LoRA subspace and that out-of-distribution loss increase is controlled by perturbation energy projected onto the leading eigenspaces of $\Sigma_T$. Then a spiked spectrum in $\Sigma_T$ is consistent with increased sensitivity to shifts aligned with that narrow adapter subspace, whereas a flatter spectrum is consistent with more distributed local robustness.
\end{corollary}

\paragraph{Significance.}
Meta-SWAG turns posterior covariance into an interpretable object: it is not just uncertainty mass, but a compressed summary of local trajectory geometry. Any deployment-robustness interpretation requires an extra assumption linking distribution shift to that geometry.

\subsection{Conceptual Limit: Self-Knowledge Bottlenecks Posterior Quality}

The next result is intentionally more interpretive than the first two. It does not characterize the implemented weighting rule directly; instead it formalizes a limitation principle for posterior predictors built from an explicit self-model $\tilde{\pi}$ of the policy being updated.

\begin{theorem}[Conceptual self-knowledge bound under an explicit self-model]
Let $\pi_\theta$ be the true policy at checkpoint $\theta$, let $\tilde{\pi}_\theta$ be the optimizer's self-model of that policy, and let $\hat{\pi}_{q_T}=\E_{\theta' \sim q_T}[\tilde{\pi}_{\theta'}]$ be the policy induced by the Meta-SWAG posterior predictive. If
\begin{equation}
C := \sup_x \left| \log \frac{\tilde{\pi}_\theta(x)}{\hat{\pi}_{q_T}(x)} \right| < \infty,
\end{equation}
then
\begin{equation}
\KL(\pi_\theta \,\|\, \hat{\pi}_{q_T})
\leq
\left(1 + \frac{C^2}{2}\right)L_{\mathrm{self}}(\theta) + O(T^{-1}),
\end{equation}
where $L_{\mathrm{self}}(\theta):=\KL(\pi_\theta \,\|\, \tilde{\pi}_\theta)$ is the self-knowledge gap.
\end{theorem}

\paragraph{Proof sketch.}
Decompose $\KL(\pi_\theta \,\|\, \hat{\pi}_{q_T})$ into $\KL(\pi_\theta \,\|\, \tilde{\pi}_\theta)$ plus a cross term involving $\log(\tilde{\pi}_\theta / \hat{\pi}_{q_T})$. Bound the cross term by a change-of-measure argument and Pinsker's inequality, then use Young's inequality to isolate the residual into a multiplicative constant on $L_{\mathrm{self}}$. The remaining $O(T^{-1})$ term is the usual finite-sample error from the SWAG moment approximation along a mixing Markovian checkpoint trajectory.

\begin{corollary}[Alignment]
Let $\pi$ be the base LLM policy, $\tilde{\pi}$ the aligner's internal model of the base policy, and $\hat{\pi}_{q_T}$ the aligned posterior predictive policy. Then alignment quality, measured as $\KL(\pi \,\|\, \hat{\pi}_{q_T})$, is upper-bounded by the aligner's self-knowledge of the base model up to the constant above. Read operationally, this means posterior averaging cannot compensate arbitrarily for a poor internal model of the policy being aligned.
\end{corollary}

\paragraph{Significance.}
The theorem identifies a limitation principle for adaptive systems: better uncertainty estimates cannot outrun a poor internal model of the policy being updated. Its practical force depends on whether one can estimate or upper-bound the self-knowledge gap in real aligners, which the current draft does not yet do for DPO, PPO, or GRPO.

\section{Experiments}

\subsection{Matrix-Game Validation}

The first empirical block validates the variance and geometry claims in controlled environments where equilibria and basin structure can be inspected directly. Our executable scaffold currently focuses on matching pennies, a 3-action zero-sum variant, and a first Kim-style iterated IPD/RPS adaptation runner with bundled personas. The comparison set is: point estimate, vanilla SWAG over checkpoints, naive softmax Meta-SWAG, and Goodhart-resilient Meta-SWAG with ESS-constrained or thresholded weighting.

For Theorem~1, the main metric is effective predictive variance of the Bayesian model average relative to both a point estimate and a uniformly weighted checkpoint ensemble. The target plot places the empirical variance ratio against the theoretical HM/AM prediction across different levels of checkpoint heterogeneity while also checking whether Goodhart-resilient weighting avoids collapse to single-checkpoint behavior. For Theorem~2, we compute the leading eigenvalues of $\Sigma_T$ and compare them to analytically or numerically estimated basin width around the convergent region. We also log the top five eigenvalues of the low-rank deviation matrix across runs to test whether the retained trajectory occupies a narrow or broad subspace.

\begin{figure}[H]
\centering
\fbox{
\begin{minipage}{0.94\linewidth}
\small
\textbf{Figure 2 (placeholder): Matrix-game validation.}

\vspace{0.4em}
\textbf{Left panel:} empirical variance ratio versus the predicted HM/AM ratio for point estimate, uniformly weighted SWAG, and Meta-SWAG variants.

\textbf{Right panel:} posterior top-eigenvalue ratio versus basin-width estimate across games and initialization regimes.
\end{minipage}
}
\caption{Main matrix-game figure for validating Theorems 1 and 2.}
\end{figure}

\subsection{LLM Alignment Validation}

The second empirical block tests Meta-SWAG on LoRA-based alignment of a Llama-family model. The draft protocol uses DPO as the headline setting and treats PPO and GRPO as algorithm-agnostic extensions rather than separate experiments \citep{rafailov2023dpo,schulman2017ppo,shao2024deepseekmath}. Checkpoints are collected along the fine-tuning trajectory; evidence scores are computed from held-out alignment metrics; the posterior is fit in LoRA space only. The main comparison is between the MAP checkpoint, standard softmax-weighted Meta-SWAG, and Goodhart-resilient Meta-SWAG under ESS-constrained or thresholded weighting.

The first target measurement is reward overoptimization. Following the best-of-$n$ evaluation style used in the alignment literature \citep{gao2023overoptimization,rafailov2024daa}, we compare the gap between a proxy reward and a gold reward model as $n$ increases. The weighting metric, model-selection metric, and final evaluation metric should be reported separately so the Goodhart-mitigation claim is not tautological. The core three conditions are: MAP checkpoint, standard softmax-weighted Meta-SWAG, and Goodhart-resilient Meta-SWAG. A Bayesian reward or Laplace-style policy baseline is included where available \citep{yang2024laplacelora,yang2024brm}. The second target measurement is algorithm-agnosticism: the same posterior construction should improve or stabilize outcomes across DPO, PPO, and GRPO without retraining ensembles for each aligner.

The present LLM protocol does not yet estimate the self-knowledge gap from Theorem~3. A faithful operationalization would require an explicit surrogate for the aligner's internal policy model, or at least a measurable proxy for mismatch between the updater's reference policy and the actual policy being aligned. Until that is implemented, Theorem~3 should be read as a limitation principle rather than as an empirically validated claim about DPO, PPO, or GRPO.

This section is intentionally placeholder-ready. The figure slots, metrics, and evaluation protocol are fixed even when some final runs are pending. Where results are not yet complete, the paper should say so explicitly rather than implying completed experiments.

\begin{figure}[H]
\centering
\fbox{
\begin{minipage}{0.94\linewidth}
\small
\textbf{Figure 3 (placeholder; some rows pending): LLM alignment validation.}

\vspace{0.4em}
\textbf{Top row:} reward-overoptimization gap versus best-of-$n$ for MAP, standard softmax Meta-SWAG, and Goodhart-resilient Meta-SWAG under DPO.

\textbf{Bottom rows:} algorithm-agnostic improvement under DPO, PPO, and GRPO measured by proxy-gold gap and posterior predictive variance.

\vspace{0.4em}
Rows without completed runs should be explicitly marked \RESULTPENDING.
\end{minipage}
}
\caption{Main LLM figure scaffold. The draft keeps this section in the body even if some numeric entries are still pending.}
\end{figure}

\paragraph{Current draft status.}
The matrix and iterated-game sections form the completed empirical base of the current draft. The LLM subsection should remain in the body with full protocol and caption text even if the numeric cells are temporarily empty.

\section{Related Work}

Meta-SWAG sits at the intersection of three threads. The first is posterior approximation in deep learning, especially SWAG, subspace inference, and adapter-space Bayesian methods \citep{maddox2019swag,yang2024laplacelora}. Our contribution is not a new low-rank posterior family; it is the claim that a SWAG-style posterior can be recovered on adaptive checkpoint processes once retention is restricted to a stable, mixing late-trajectory regime.

The second thread is meta-learning in multi-agent reinforcement learning. Kim-style analyses, LOLA-style opponent shaping, and related methods describe the dynamics of adaptation, but they do not typically expose a cheap posterior over checkpoints that can be sampled at deployment time \citep{kim2021meta,foerster2018lola,giannou2022convergence}. Meta-SWAG is intended as a posterior layer that can sit on top of adaptive optimization trajectories, provided checkpoint retention and held-out evidence are available.

The third thread is Bayesian alignment. Reward-model posteriors and Laplace-style approximations quantify uncertainty over evaluators or local parameter neighborhoods \citep{yang2024laplacelora,yang2024brm}. Meta-SWAG instead quantifies uncertainty over the aligned policy trajectory and is therefore closest in spirit to policy-side Bayesian model averaging.

\section{Discussion and Limitations}

Meta-SWAG is deliberately modest about what it approximates. The posterior is Gaussian, which is unlikely to be adequate when alignment basins are genuinely bimodal or when policy trajectories move between distant equilibria. Markovianity alone is not enough to justify the approximation: the argument also needs retention inside a stable, mixing region where late-trajectory moments are informative. The covariance interpretation is strongest near convergence, where the local second-order story is meaningful; earlier in training, the low-rank factor is better interpreted as trajectory structure than as calibrated epistemic uncertainty.

The self-knowledge result is also conservative and should be read as conceptual unless an explicit self-model is implemented in the empirical stack. The multiplicative constant $1 + C^2/2$ is easy to derive and easy to explain, but it is unlikely to be tight for expressive policies. Richer divergences or chi-squared style refinements may sharpen the bound, but the current form already isolates the main conceptual point: posterior quality in adaptive systems is bottlenecked by self-knowledge, not just by moment estimation.

The most natural extension is to richer posterior families such as MultiSWAG or mixture approximations over multiple alignment basins. A second direction is to combine Meta-SWAG with causal or mechanistic diagnostics so that posterior samples are not only uncertain but also interpretable. For the present paper, however, the goal is narrower: establish that one retained-checkpoint posterior construction can be used on adaptive learning trajectories in both MARL and LLM alignment, without claiming that the two domains share the same underlying non-stationarity beyond that abstraction.

\appendix

\section{Notation and Assumptions}

We use $\theta_t$ for checkpoints in both application domains.
\begin{itemize}[leftmargin=1.5em]
  \item In MARL, $\theta_t$ denotes the joint policy parameters at meta-time $t$.
  \item In LLM alignment, $\theta_t$ denotes the LoRA adapter parameters at step $t$; the base model remains frozen.
  \item $m_t$ is an evidence score on a held-out meta-objective.
  \item $w_t$ is the normalized evidence weight derived from $m_t$.
  \item $q_T(\theta)$ is the Meta-SWAG posterior after $T$ retained checkpoints.
\end{itemize}

The draft assumes: (i) Markovian trajectory coherence; (ii) late-trajectory concentration in a region where moment estimation is meaningful; (iii) bounded moments; and (iv) availability of a held-out evidence score for each retained checkpoint. The Gaussian approximation is justified by the combination of (i)--(iii), not by Markovianity alone.

\section{Posterior Construction Details}

\subsection{Evidence Weighting}

The draft main text treats naive softmax weighting as a baseline
\begin{equation}
w_t = \frac{\exp(\beta m_t)}{\sum_{s=1}^T \exp(\beta m_s)}.
\end{equation}
This rule is scale-free once $m_t$ values are centered, but by itself it is vulnerable to Goodhart-style posterior collapse when $m_t$ is a proxy objective. The current implementation therefore supports two Goodhart-resilient alternatives: ESS-constrained softmax weighting, which adapts $\beta$ to preserve a minimum effective sample size, and thresholded satisficing weights, which assign uniform mass to all checkpoints above a validation threshold.

\subsection{Diagonal-Plus-Low-Rank Covariance}

The covariance estimator is
\begin{equation}
\Sigma_T = \frac{1}{2}\diag(v_T) + \frac{1}{2(|\mathcal{K}|-1)}D_TD_T^\top.
\end{equation}
The diagonal term controls per-parameter uncertainty; the low-rank term preserves the dominant trajectory subspace. In the LLM setting the computation is performed only in LoRA space, which is critical both computationally and conceptually.

\subsection{Posterior Prediction}

Given samples $\theta^{(s)} \sim q_T$, posterior prediction uses
\begin{equation}
p(y \mid x, \mathcal{D}) \approx \frac{1}{S}\sum_{s=1}^S p(y \mid x, \theta^{(s)}).
\end{equation}
For MARL, the same averaging is performed over policy outputs or rollout-level statistics. For LLMs, the default deployment object is a Bayesian model average over LoRA-sampled policies rather than a single merged checkpoint.

\section{Expanded Theory Notes}

\subsection{Theorem 1 Notes}

The first theorem should be read as a posterior analogue of the variance reduction seen in evidence-weighted policy gradients. A longer proof path for the paper version is: decompose predictive variance into within-checkpoint and between-checkpoint components; show that evidence weighting controls the within-checkpoint component by the harmonic mean of checkpoint-level variances; and show that the between-checkpoint term vanishes in the degenerate point-estimate limit and remains controlled whenever the retained trajectory lies in a single basin. The reward-overoptimization corollary is not automatic from this argument. It needs an extra bridge assumption tying proxy-gold gap to predictive dispersion on the retained posterior support.

\subsection{Theorem 2 Notes}

The geometry theorem depends on a local stationary approximation. In the final write-up it will help to isolate two claims that are currently bundled together: the retained trajectory mixes within a stable basin, and the empirical covariance estimates the local second-order geometry of that basin. The theorem does not claim global posterior calibration, and any robustness-under-shift interpretation requires a separate assumption linking deployment shifts to the retained LoRA subspace.

\subsection{Theorem 3 Notes}

The self-knowledge theorem is a limitation principle rather than a completed empirical theorem for the current stack, so the appendix should carry a slightly more explicit derivation together with a short note about what would be required to instantiate $L_{\mathrm{self}}$ for DPO, PPO, or GRPO. Start from
\begin{equation}
\KL(\pi_\theta \,\|\, \hat{\pi}_{q_T})
=
\KL(\pi_\theta \,\|\, \tilde{\pi}_\theta)
+ \E_{\pi_\theta}\left[\log \frac{\tilde{\pi}_\theta}{\hat{\pi}_{q_T}}\right].
\end{equation}
Then identify the first term with $L_{\mathrm{self}}(\theta)$, bound the second term using the uniform log-density ratio constant $C$, move from total variation to KL using Pinsker, absorb the mixed term with Young's inequality, and append the finite-sample posterior approximation term $O(T^{-1})$.

\section{Experiment Protocols}

\subsection{Matrix Games}

Planned environments are matching pennies, a 3-action zero-sum variant, and Kim-style iterated IPD/RPS follow-ups. Planned comparisons are point estimate, vanilla SWAG over retained checkpoints, standard softmax Meta-SWAG, and Goodhart-resilient Meta-SWAG with ESS-constrained or thresholded weighting. Primary metrics are effective predictive variance, HM/AM-predicted versus empirical variance ratio, effective sample size, top-eigenvalue concentration of $\Sigma_T$, and a basin-width proxy from convergence under perturbed initialization.

\subsection{LLM Alignment}

Planned setup: a Llama-family base model, LoRA-only adaptation, DPO as the primary aligner with PPO and GRPO extensions, an Alpaca-style preference benchmark or equivalent, and evaluation through proxy reward, gold reward model, best-of-$n$ gap, posterior predictive variance, effective sample size, and max-weight collapse diagnostics. If results are not finished, the paper should still report the checkpoint collection rule, the evidence score definition, the evaluation metrics, the intended figure captions, and which rows are pending.

\section{LoRA Tracking Implementation Notes}

The implementation should track only the trainable LoRA tensors during alignment. A practical first pass is: save LoRA adapter weights every fixed number of optimizer steps after burn-in; score each retained checkpoint on a held-out alignment split; compute evidence weights from the held-out score; accumulate the weighted mean and diagonal moments online; retain only the latest $K$ centered deviations for the low-rank factor; and sample LoRA adapters at evaluation time before merging them into the frozen base model.

\bibliographystyle{plainnat}
\bibliography{references}

\end{document}
